\subsection{Jogos em grafos arbitrários}

Suponha que dois jogadores jogam um jogo num grafo arbitrário $G$. Os jogadores se movem em turnos e vão para um vértice adjacente. Dependendo do jogo, a pessoa que não puder se mover vai perder ou ganhar.

Consideramos um grafo dirigido arbitrário. A tarefa é determinado, dado um estado inicial, quem vai ganhar o jogo se os dois tiverem estratégias ótimas (ou então determinar que haverá um empate). Será encontrada uma solução para todos os vértices iniciais do grafo em $O(m)$ onde $m=$ num de arestas.

Um vértice é ganhador se um player começando nele ganha o jogo se jogar de maneira ótima. Analogamente denotaremos o outro por perdedor. Para todos os vértices sem outgoing edges já sabemos sua condição de vitória/derrota. Além disso:

\begin{enumerate}
    \item Se um vértice tem uma aresta que vai para um perdedor, então ele é ganhador
    \item Se todas as arestas que saem de um vértice são vencedoras, então o vértice é perdedor
    \item Se em algum ponto ainda existirem vértices indefinidios, então é um vértice de empate
\end{enumerate}

Então dá pra fazer um algoritmo em $O(nm)$: roda pra todos os vértices aplicando essas regras. Mas da pra reduzir pra $O(m)$.

Iteramos por todos os vértices que sabemos o estado inicial. Para cada um deles fazemos uma DFS que se move considerando as arestas ao contrário. Ela não vai entrar em vértices que já tem vitória/derrota definida. Se a pesquisa for de um vértice perdedor para um indefinido, então marca o indefinido como ganhador e continua a DFS dele. Se nós formos de um ganhador para um indefinido, então chegamos se todas as arestas desse indefinido vão para ganhadores. Dá pra fazer isso em $O(1)$ guardando o número de arestas que vão para vencedores para cada vértice.

\inccode{games/generic_graph_game.cpp}

\prob{1}{gamesongraphs-1} Polícia e ladrão. Tem um tabuleiro $m \times n$. Algumas das células não podem ser entradas. As coordenadas iniciais do policial e do ladrão são conhecidas. Uma das células é a saída. Se o policial e o ladrão estão na mesma cela em qualquer momento, o policial ganha. Se o ladrão estiver na saída (sem que o policial esteja lá também), ele ganha. O policial pode andar nas $8$ direções e o ladrão em $4$. O policial e o ladrão alternam os movimentos, mas também podem skipar um movimento se quiserem.

\textbf{Solução:} construímos um grafo. O estado do jogo é definido pelas coordenadas do policial, do ladrão e de quem é a vez. Ou seja, um vértice do grafo é $(P, T, P_turn)$. Precisamos determinar quais vértices são perdedores e ganhadores. Se é a vez do policial, então o vértice é um winning se as coordenadas dos dois coincidirem e é losing se não for winning E o policial está na exit. Se é a vez do ladrão, então é losing se as coordenadas coincidirem e winning se não for losing e o ladrão estiver no exit vertex.

Sobre a implementação, é importante decidir se você vai construir o grafo explicitamente ou on the fly. O pro é que tem menos chance de ter erro e vai ser bem mais fácil de codar. O contra é que vai ser muito lento do que usando on the fly.

Aqui uma approach construindo explicitamente:

\inccode{games/policia_e_ladrao.cpp}